\chapter*{Hogyan használd ezt a könyvet}
Minden feladat ugyanúgy épül fel, hogy vizsgán is ugyanebben a rendben tudj gondolkodni.
\begin{enumerate}
\item \textbf{Feladat.} A feladat szövege, alatta a forrása (vizsga, gyűjtemény, oldal).
\item \textbf{Mit keresünk?} Minden részkérdésnél leírom, melyik mennyiséget keressük, és mit jelent.
\item \textbf{Jelölések ebben a feladatban.} Minden jel jelentése abban a környezetben, ahol épp használjuk, és ha a feladat megadja, az értéke is. Ugyanaz a betű a tárgyban több dolgot is jelenthet (a $k$ lehet Nyquist-zóna sorszáma, felharmonikus sorszáma vagy erősítési tényező), ezért ezt minden feladatnál újra rögzítem.
\item \textbf{Megoldás lépésenként.} Minden matematikai sor alatt kis kapitálissal áll a sor jellege (\textsc{tétel}, \textsc{behelyettesítés}, \textsc{számérték}, \ldots) és egy rövid megjegyzés arról, honnan jön.
\item \textbf{Eredmény}, és szükség esetén \textbf{Figyelem} doboz a tipikus hibákkal, illetve a forrásban talált tévedésekkel (\textbf{Javítás}).
\end{enumerate}

\section*{A színek jelentése}
A könyv fekete-fehér, egyedül a behelyettesítést jelzem színekkel. A feladat által megadott mennyiségek közül legfeljebb ötnek saját színe van: a jelöléstáblában, a képletben a jele, a behelyettesítésnél az értéke ugyanabban a színben jelenik meg. Így ránézésre látszik, mi hova kerül.
\begin{center}
\ca{kék}\quad \cb{narancs}\quad \cc{lila}\quad \cd{kékeszöld}\quad \ce{bíbor}
\end{center}
\behely{\bh{\ca{f_H}}{\ca{75\un{kHz}}}}
\[
f_s\ge2\,\ca{f_H}=2\cdot\ca{75}=150\un{kHz}
\]

\section*{Az ábrákról}
Ahol a feladat rajzot kér, a \textbf{Rajz: ezt kell felrajzolni} doboz pontos ábrát ad, alatta pontokba szedve, hogy kézzel mit kell feltüntetni (tengelyek, jellegzetes értékek, burkoló). Minden ábra papíron egyetlen tollal is megrajzolható.

\section*{Ellenőrzés}
Minden számértéket független szkript számolt újra (\texttt{html/check}). A források hibáit a Javítás doboz és a hibajegyzék (\texttt{03\_Hibajegyzek.md}) rögzíti.
